\documentclass{amsart}
\usepackage[english]{babel}
\usepackage{geometry,amssymb,hyperref}
\geometry{letterpaper}

%%%%%%%%%% Start TeXmacs macros
\newcommand{\assign}{:=}
\newcommand{\cdummy}{\cdot}
\newcommand{\nin}{\not\in}
\newcommand{\tmSep}{; }
\newcommand{\tmem}[1]{{\em #1\/}}
\newcommand{\tmmathbf}[1]{\ensuremath{\boldsymbol{#1}}}
\newcommand{\tmop}[1]{\ensuremath{\operatorname{#1}}}
\newcommand{\tmsep}{, }
\newcommand{\tmtextbf}[1]{\text{{\bfseries{#1}}}}
\newcommand{\tmtextit}[1]{\text{{\itshape{#1}}}}
\newcommand{\tmtextup}[1]{\text{{\upshape{#1}}}}
\newtheorem{corollary}{Corollary}
\newtheorem{definition}{Definition}
\newtheorem{lemma}{Lemma}
\newtheorem{proposition}{Proposition}
{\theoremstyle{remark}\newtheorem{remark}{Remark}}
\newtheorem{theorem}{Theorem}
%%%%%%%%%% End TeXmacs macros

\newcommand{\Ima}{\tmop{Im}}
\newcommand{\Rea}{\tmop{Re}}
\newcommand{\cum}{cum}
\newcommand{\supp}{\tmop{supp}}
\newcommand{\abs}[1]{\left| #1 \right|}
\newcommand{\norm}[1]{|#1|}
\newcommand{\Zt}{\tilde{Z}}

\begin{document}

\title{Band-Limitedness of the Hardy $Z$-Function\\
in the Phase Variable\\
and its Higher Cumulant Spectra}

\author{Stephen Crowley}
\email{stephencrowley214@gmail.com}

\date{April 29, 2026}

\begin{abstract}
  Let $\vartheta (t) = \Ima \log \Gamma \hspace{-0.17em} \left( \tfrac{1}{4} +
  \tfrac{it}{2} \right) - \tfrac{t}{2} \log \pi$ be the Riemann--Siegel theta
  function and $Z (t) = e^{i \vartheta (t)} \zeta \hspace{-0.17em} \left(
  \tfrac{1}{2} + it \right)$ the Hardy $Z$-function. Fix any constant $C >
  c_{\ast}$, where
  \[ c_{\ast} \hspace{0.27em} : = \hspace{0.27em} \tfrac{1}{2}
     \hspace{-0.17em} \left( \gamma + 3 \log 2 + \tfrac{\pi}{2} + \log \pi
     \right) \hspace{0.27em} \approx \hspace{0.27em} 2.6860917, \]
  so that $\Theta (t) : = \vartheta (t) + Ct$ is a $C^{\infty}$ diffeomorphism
  of $[0, \infty)$, and let $\Zt : = Z \circ \Theta^{- 1}$. We prove that for
  every $\xi \in \mathbb{R}$ with $\abs{\xi} > 1$,
  \[ \frac{1}{\Theta (T)} \int_0^{\Theta (T)} \Zt (u) \hspace{0.17em} e^{- i
     \xi u} \hspace{0.17em} du \hspace{0.27em} \longrightarrow \hspace{0.27em}
     0 \quad \text{as} T \to \infty, \]
  so the Ces{\`a}ro spectral support of $\Zt$ is contained in $[- 1, 1]$. The
  proof rests on three ingredients: (i)~the Riemann--Siegel expansion with
  remainder $O (t^{- 1 / 4}$); (ii)~a single integration by parts whose tail
  is eliminated by an exact antiderivative; and (iii)~the elementary bound
  $\omega_n (u) \in [0, 1$), where $u_n$ is the exact saddle of $\omega_n$.
  The cutoff $\abs{\xi} = 1$ is sharp. We then extend the argument to
  higher-order cumulant spectra: for every $s \ge 2$ the $s$-point
  polyspectrum $\Phi_s$ of $\Zt$ is supported in the box $[- 1, 1]^{s - 1}$.
\end{abstract}

\subjclass{11M06\tmSep  11M26\tmSep  42A38\tmSep  60G10\tmSep  62M15}

\keywords{Hardy $Z$-function\tmsep  Riemann--Siegel formula\tmsep 
band-limited function\tmsep  phase variable\tmsep  Ces{\`a}ro--Fourier
transform\tmsep  polyspectrum\tmsep  cumulant}

{\maketitle}

{\tableofcontents}

\section{Introduction}

The Riemann--Siegel theta function $\vartheta (t)$ satisfies $\vartheta' (t)
\sim \tfrac{1}{2} \log (t / 2 \pi) \to + \infty$, so the Hardy $Z$-function
oscillates at a rate that grows without bound. A natural strategy is to absorb
this growth into a change of variables: replace the time parameter $t$ by the
{\tmem{accumulated phase}} $u = \Theta (t) \assign \vartheta (t) + Ct$, where
$C > c_{\ast}$ is chosen so that $\Theta$ is a globally increasing
diffeomorphism. The reparametrized function $\Zt (u) \assign Z (\Theta^{- 1}
(u))$ then admits a well-defined spectral theory in the Ces{\`a}ro sense.

The main result, Theorem~\ref{thm:band}, shows that the Ces{\`a}ro--Fourier
transform of $\Zt$ is supported in $[- 1, 1]$. The argument proceeds in three
steps that mirror the Riemann--Siegel decomposition~\eqref{eq:RS}. The cosine
sum contributes pairs of mode integrals $I_n (\xi, T)$ and $J_n (\xi, T)$,
each bounded via a single integration by parts
(Lemmas~\ref{lem:IBP}--\ref{lem:IBP-companion}); the error term is absorbed by
the classical $O (t^{- 1 / 4})$ bound on the Riemann--Siegel remainder.

A distinguishing feature of this argument is that the tail integral in the
integration by parts telescopes {\tmem{exactly}}, yielding the sharp bound
$\abs{I_n (\xi, T)} \le 2 / (\xi - 1)$ with no asymptotic error. The cutoff
$\abs{\xi} = 1$ is confirmed to be optimal in Corollary~\ref{cor:sharp}: at
every interior frequency $\xi \in (- 1, 1)$ the phase has a non-degenerate
stationary point, and after a mode-scale separation argument the Ces{\`a}ro
average remains bounded below.

{\noindent}\tmtextbf{Higher cumulant spectra.} Section~\ref{sec:higher}
extends these results to the $s$-point {\tmem{polyspectrum}} $\Phi_s$ of
$\Zt$, i.e., the Fourier transform of the $s$-point Ces{\`a}ro cumulant
density $\kappa_s$. We show that $\Phi_s$ is supported in $[- 1, 1]^{s - 1}$
(Theorem~\ref{thm:poly}), generalizing the band-limitedness of the two-point
case. The proof uses a multi-mode IBP in which each frequency variable $\xi_j$
is handled independently via Lemma~\ref{lem:IBP}.

{\noindent}\tmtextbf{Corrections to a prior draft.} This revision fixes three
issues in an earlier version of the manuscript {\cite{bandlim-src}}. (i)~The
identity $\vartheta'  (2 \pi n^2) = \log n$ is only asymptotic: $\vartheta' 
(2 \pi n^2) = \log n - 1 / (16 \pi^2 n^4) + O (n^{- 8})$. We therefore
redefine $u_n$ as $\Theta (t_n^{\dagger})$, where $t_n^{\dagger}$ is the
{\tmem{exact}} solution of $\vartheta' (t) = \log n$, so that $\omega_n (u_n)
= 0$ is an identity by construction. (ii)~The identity $2 \cos \Phi_n = e^{i
\Phi_n} + e^{- i \Phi_n}$ produces two distinct mode integrals $I_n$ and
$J_n$; the claim ``$2 \hspace{0.17em} \Rea I_n$'' in the earlier draft
conflates them. We bound both by separate (identical) IBPs. (iii)~The
mode-separation step in the sharpness proof is made explicit using the
asymptotic scales $t_n^{\ast} \sim 2 \pi n^{2 / (1 - \xi)} e^{2 \xi C / (1 -
\xi)}$.

{\noindent}\tmtextbf{Notation.} $\gamma$ denotes the Euler--Mascheroni
constant, $\psi = \Gamma' / \Gamma$ the digamma function, $f \ll g$ means
$\abs{f} \le Cg$ for an absolute constant $C > 0$, and $\lfloor
\hspace{0.17em} \cdummy \hspace{0.17em} \rfloor$ is the floor function.

\section{Setup and Notation}

\begin{definition}
  [Hardy $Z$-function]\label{def:Z} For $t \ge 0$ define
  \[ \vartheta (t) \hspace{0.27em} \assign \hspace{0.27em} \Ima \log \Gamma
     \hspace{-0.17em} \left( \tfrac{1}{4} + \tfrac{it}{2} \right) -
     \tfrac{t}{2} \log \pi, \qquad Z (t) \hspace{0.27em} \assign
     \hspace{0.27em} e^{i \vartheta (t)} \zeta \hspace{-0.17em} \left(
     \tfrac{1}{2} + it \right) . \]
  The function $Z$ is real-valued with zeros exactly the zeros of $\zeta$ on
  $\Rea s = \tfrac{1}{2}$.
\end{definition}

\begin{definition}
  [Shifted phase map]\label{def:Theta} Set
  \[ c_{\ast} \hspace{0.27em} \assign \hspace{0.27em} \tfrac{1}{2}
     \hspace{-0.17em} \left( \gamma + 3 \log 2 + \tfrac{\pi}{2} + \log \pi
     \right), \]
  fix any $C > c_{\ast}$, and define
  \[ \Theta (t) \hspace{0.27em} \assign \hspace{0.27em} \vartheta (t) + Ct,
     \qquad g \hspace{0.27em} \assign \hspace{0.27em} \Theta^{- 1}, \qquad \Zt
     (u) \hspace{0.27em} \assign \hspace{0.27em} Z (g (u)) . \]
\end{definition}

\begin{proposition}
  [Diffeomorphism property]\label{prop:diffeo} For any $C > c_{\ast}$, the map
  $\Theta : [0, \infty) \to [\Theta (0), \infty)$ is a $C^{\infty}$
  diffeomorphism with
  \[ \Theta' (t) \hspace{0.27em} = \hspace{0.27em} \vartheta' (t) + C
     \hspace{0.27em} > \hspace{0.27em} 0 \quad \forall \hspace{0.17em} t \ge
     0, \qquad \Theta' (t) \hspace{0.27em} \sim \hspace{0.27em} \tfrac{1}{2}
     \log \hspace{-0.17em} \left( \tfrac{t}{2 \pi} \right)  \quad (t \to
     \infty) . \]
  Moreover, $\Theta (T) \sim \tfrac{T}{2} \log (T / 2 \pi)$ as $T \to \infty$.
\end{proposition}

\begin{proof}
  The identity $\vartheta' (t) = \tfrac{1}{2} \Rea \hspace{0.17em} \psi
  \hspace{-0.17em} \left( \tfrac{1}{4} + \tfrac{it}{2} \right) - \tfrac{1}{2}
  \log \pi$ and the digamma special value $\psi \hspace{-0.17em} \left(
  \tfrac{1}{4} \right) = - \gamma - \tfrac{\pi}{2} - 3 \log 2$ give
  \[ \vartheta' (0) = \tfrac{1}{2} \psi \hspace{-0.17em} \left( \tfrac{1}{4}
     \right) - \tfrac{1}{2} \log \pi = - \tfrac{1}{2} \hspace{-0.17em} \left(
     \gamma + \tfrac{\pi}{2} + 3 \log 2 + \log \pi \right) = - c_{\ast} . \]
  Since $\vartheta'$ is strictly increasing with $\vartheta' (t) =
  \tfrac{1}{2} \log (t / 2 \pi) + O (t^{- 2})$ {\cite{Titchmarsh1986}}, its
  global infimum is $- c_{\ast}$, so $\Theta' (t) = \vartheta' (t) + C > 0$
  for all $t \ge 0$ whenever $C > c_{\ast}$. Surjectivity onto $[\Theta (0),
  \infty)$ is immediate from $\Theta (t) \to \infty$. The asymptotic for
  $\Theta (T)$ follows by integration.
\end{proof}

\section{Riemann--Siegel Expansion}

Set $N (t) \assign \lfloor \sqrt{t / (2 \pi)} \rfloor$. The Riemann--Siegel
formula {\cite{Siegel1932,Titchmarsh1986}} asserts
\begin{equation}
  \label{eq:RS} Z (t) \hspace{0.27em} = \hspace{0.27em} 2 \sum_{n = 1}^{N (t)}
  n^{- 1 / 2} \cos \hspace{-0.17em} (\vartheta (t) - t \log n) + R (t), \qquad
  \abs{R (t)} \ll t^{- 1 / 4}, \quad t \ge 1.
\end{equation}
\begin{lemma}
  [Partial-sum bound]\label{lem:partial} For all $T \ge 1$,
  \[ \sum_{n = 1}^{N (T)} n^{- 1 / 2} \hspace{0.27em} \le \hspace{0.27em} 2
     \hspace{-0.17em} \left( \frac{T}{2 \pi} \right)^{1 / 4} + 2. \]
\end{lemma}

\begin{proof}
  The standard estimate $\sum_{n = 1}^M n^{- 1 / 2} \le 2 \sqrt{M} + 1$,
  applied with $M = N (T) \le (T / 2 \pi)^{1 / 2}$, gives the result.
\end{proof}

\section{Phase-Derivative Identity}

\begin{definition}
  [Exact saddle points]\label{def:saddle} For each $n \ge 1$ let
  $t_n^{\dagger}$ denote the unique solution in $(0, \infty)$ of
  \[ \vartheta' (t_n^{\dagger}) = \log n, \]
  and set $u_n \assign \Theta (t_n^{\dagger})$. By the asymptotic $\vartheta'
  (t) = \tfrac{1}{2} \log (t / 2 \pi) + O (t^{- 2})$,
  \[ t_n^{\dagger} \hspace{0.27em} = \hspace{0.27em} 2 \pi n^2 + O (n^{- 2}),
     \qquad u_n \hspace{0.27em} = \hspace{0.27em} \Theta (2 \pi n^2) + O
     \hspace{-0.17em} \left( \frac{\log n}{n^2} \right) . \]
\end{definition}

\begin{remark}
  An earlier draft {\cite{bandlim-src}} used $u_n = \Theta (2 \pi n^2)$
  directly, relying on the {\tmem{asymptotic}} identity $\vartheta'  (2 \pi
  n^2) = \log n$. The correction is of order $n^{- 4}$: $\vartheta'  (2 \pi
  n^2) = \log n - \frac{1}{16 \pi^2 n^4} + O (n^{- 8})$. Using $t_n^{\dagger}$
  in place of $2 \pi n^2$ makes the saddle condition $\omega_n (u_n) = 0$ an
  exact identity, not an asymptotic one, at the negligible cost of an $O (n^{-
  4})$ perturbation of the lower limit of the mode integral.
\end{remark}

For $n \ge 1$ and $u \ge u_n$, define the {\tmem{mode phase}} and
{\tmem{frequency ratio}}
\begin{equation}
  \label{eq:Phi-omega} \Phi_n (u) \hspace{0.27em} \assign \hspace{0.27em}
  \vartheta (g (u)) - g (u) \log n, \qquad \omega_n (u) \hspace{0.27em}
  \assign \hspace{0.27em} \frac{\vartheta' (g (u)) - \log n}{\Theta' (g (u))}
  .
\end{equation}
\begin{lemma}
  [Phase-derivative identity]\label{lem:phase} For all integers $n \ge 1$, all
  $u \ge u_n$, and all $\xi \in \mathbb{R}$:
  \begin{enumerate}
    \item[\tmtextup{(a)}] $\frac{d}{du}  (\Phi_n (u) - \xi u) = \omega_n (u) -
    \xi$.
    
    \item[\tmtextup{(b)}] $\omega_n (u_n) = 0$ {\tmem{exactly}}, by
    Definition~\ref{def:saddle}.
    
    \item[\tmtextup{(c)}] $0 \le \omega_n (u) < 1$ for all $u \ge u_n$.
    
    \item[\tmtextup{(d)}] $\omega_n' (u) = \frac{(C + \log n)  \hspace{0.17em}
    \vartheta'' (g (u))}{\Theta' (g (u))^3} > 0$; hence $\omega_n$ is strictly
    increasing on $[u_n, \infty)$.
    
    \item[\tmtextup{(e)}] For every constant $\xi \nin \omega_n
    \hspace{-0.17em} ([u_n, \infty))$,
    \[ \frac{d}{du} \hspace{-0.17em} \left[ \frac{1}{\xi - \omega_n (u)}
       \right] = \frac{\omega_n' (u)}{(\xi - \omega_n (u))^2} . \]
  \end{enumerate}
\end{lemma}

\begin{proof}
  \tmtextbf{(a).} Using $g' (u) = 1 / \Theta' (g (u))$:
  \[ \frac{d \Phi_n}{du} = (\vartheta' (g (u)) - \log n) \cdot g' (u) =
     \frac{\vartheta' (g (u)) - \log n}{\Theta' (g (u))} = \omega_n (u) . \]
  Subtracting $\xi$ gives part~(a).
  
  {\noindent}\tmtextbf{(b).} By Definition~\ref{def:saddle}, $g (u_n) =
  t_n^{\dagger}$ and $\vartheta' (t_n^{\dagger}) = \log n$, so $\omega_n (u_n)
  = 0 / \Theta' (t_n^{\dagger}) = 0$.
  
  {\noindent}\tmtextbf{(c).} Non-negativity: $\omega_n (u_n) = 0$ and
  $\omega_n$ is increasing by~(d), so $\omega_n (u) \ge 0$. The upper bound
  follows from the algebraic identity
  \[ \omega_n (u) = 1 - \frac{\log n + C}{\vartheta' (g (u)) + C} < 1, \]
  since $\log n + C > 0$ for all $n \ge 1$ and $C > 0$.
  
  {\noindent}\tmtextbf{(d).} Write $\omega_n = (f \circ g) / (h \circ g)$ with
  $f (t) \assign \vartheta' (t) - \log n$ and $h (t) \assign \Theta' (t) =
  \vartheta' (t) + C$. Then $f' = h' = \vartheta''$, and the quotient rule
  with $g' = 1 / (h \circ g)$ gives
  
  \begin{align*}
    \omega_n' (u) & = \frac{f' (g)  \hspace{0.17em} h (g) - f (g) 
    \hspace{0.17em} h' (g)}{h (g)^2} \cdot g' (u)\\
    & = \frac{\vartheta'' (g)  [h (g) - f (g)]}{h (g)^3} = \frac{(C + \log n)
    \hspace{0.17em} \vartheta'' (g (u))}{\Theta' (g (u))^3} .
  \end{align*}
  
  Positivity $\vartheta'' (t) > 0$ for all $t \ge 0$ is classical
  {\cite[5.15.3]{DLMF}}.
  
  {\noindent}\tmtextbf{(e).} Direct computation: $\tfrac{d}{du} (\xi -
  \omega_n)^{- 1} = \omega_n'  (\xi - \omega_n)^{- 2}$.
\end{proof}

\section{The Uniform IBP Bound}

For $\xi > 1$, $n \ge 1$, and $T$ with $\Theta (T) > u_n$, define the
{\tmem{mode integrals}}
\begin{equation}
  \label{eq:In-Jn} I_n (\xi, T) \assign \int_{u_n}^{\Theta (T)} e^{i (\Phi_n
  (u) - \xi u)}  \hspace{0.17em} du, \qquad J_n (\xi, T) \assign
  \int_{u_n}^{\Theta (T)} e^{- i (\Phi_n (u) + \xi u)}  \hspace{0.17em} du.
\end{equation}
\begin{lemma}
  [Uniform IBP bound for $I_n$]\label{lem:IBP} For every $\xi > 1$, $n \ge 1$,
  and $T$ with $\Theta (T) > u_n$,
  \[ \abs{I_n (\xi, T)} \hspace{0.27em} \le \hspace{0.27em} \frac{2}{\xi - 1}
     . \]
\end{lemma}

\begin{proof}
  Set $A \assign \xi - \omega_n (\Theta (T))$ and $B \assign \xi - \omega_n
  (u_n) = \xi$. By Lemma~\ref{lem:phase}(c), $A \ge \xi - 1 > 0$ and $B = \xi
  > 1$.
  
  {\noindent}\tmtextit{Integration by parts.} Write the integrand as $e^{i
  (\Phi_n - \xi u)} = (i (\omega_n - \xi))^{- 1} \hspace{0.17em} \tfrac{d}{du}
  e^{i (\Phi_n - \xi u)}$ and integrate by parts once:
  \begin{equation}
    \label{eq:IBP-formula} I_n (\xi, T) = \left[ \frac{e^{i (\Phi_n - \xi
    u)}}{i (\omega_n - \xi)} \right]_{u_n}^{\Theta (T)} + \int_{u_n}^{\Theta
    (T)} \frac{e^{i (\Phi_n - \xi u)}  \hspace{0.17em} \omega_n' (u)}{i (\xi -
    \omega_n (u))^2}  \hspace{0.17em} du.
  \end{equation}
  {\noindent}\tmtextit{Boundary terms.} The upper boundary contributes at most
  $1 / A$; the lower contributes $1 / B$.
  
  {\noindent}\tmtextit{Tail integral.} Since $\omega_n' > 0$ by
  Lemma~\ref{lem:phase}(d), the modulus of the tail integrand is non-negative,
  and Lemma~\ref{lem:phase}(e) supplies the exact antiderivative:
  \[ \int_{u_n}^{\Theta (T)} \frac{\omega_n' (u)}{(\xi - \omega_n (u))^2} 
     \hspace{0.17em} du \hspace{0.27em} = \hspace{0.27em} \left[ \frac{1}{\xi
     - \omega_n (u)} \right]_{u_n}^{\Theta (T)} = \frac{1}{A} - \frac{1}{B} .
  \]
  This is an identity, not an asymptotic.
  
  {\noindent}\tmtextit{Combining.}
  \[ \abs{I_n (\xi, T)} \hspace{0.27em} \le \hspace{0.27em} \frac{1}{A} +
     \frac{1}{B} + \hspace{-0.17em} \left( \frac{1}{A} - \frac{1}{B} \right)
     \hspace{0.27em} = \hspace{0.27em} \frac{2}{A} \hspace{0.27em} \le
     \hspace{0.27em} \frac{2}{\xi - 1} . \]
\end{proof}

\begin{lemma}
  [Uniform IBP bound for $J_n$]\label{lem:IBP-companion} For every $\xi > 1$,
  $n \ge 1$, and $T$ with $\Theta (T) > u_n$,
  \[ \abs{J_n (\xi, T)} \hspace{0.27em} \le \hspace{0.27em} \frac{2}{\xi}
     \hspace{0.27em} \le \hspace{0.27em} \frac{2}{\xi - 1} . \]
\end{lemma}

\begin{proof}
  The phase of the $J_n$ integrand is $- (\Phi_n (u) + \xi u)$, with
  derivative $- (\omega_n (u) + \xi)$. By Lemma~\ref{lem:phase}(c), $\omega_n
  (u) + \xi \ge \xi > 1$. Apply the same IBP as in Lemma~\ref{lem:IBP}, with
  $\xi - \omega_n$ replaced by $\xi + \omega_n$. Set $A' \assign \xi +
  \omega_n (\Theta (T))$ and $B' \assign \xi + \omega_n (u_n) = \xi$. Both $A'
  \ge \xi$ and $B' = \xi$. The boundary contributes $1 / A' + 1 / B' \le 2 /
  \xi$, and the tail antiderivative $- [(\xi + \omega_n)^{- 1}]_{u_n}^{\Theta
  (T)} = 1 / B' - 1 / A'$ contributes at most $1 / \xi - 1 / A'$. Combining:
  \[ \abs{J_n (\xi, T)} \hspace{0.27em} \le \hspace{0.27em} \frac{1}{A'} +
     \frac{1}{B'} + \hspace{-0.17em} \left( \frac{1}{B'} - \frac{1}{A'}
     \right) \hspace{0.27em} = \hspace{0.27em} \frac{2}{B'} \hspace{0.27em} =
     \hspace{0.27em} \frac{2}{\xi} . \]
\end{proof}

\begin{remark}
  In both lemmas the $1 / B$ (resp. $1 / A'$) contributions cancel exactly
  between a boundary term and the tail antiderivative. No oscillatory phase
  cancellation is invoked.
\end{remark}

\section{Band-Limitedness}

\begin{theorem}
  [Band-limitedness of $\Zt$]\label{thm:band} Let $C > c_{\ast}$, $g =
  \Theta^{- 1}$, and $\Zt = Z \circ g$. For every $\xi \in \mathbb{R}$ with
  $\abs{\xi} > 1$,
  \[ \frac{1}{\Theta (T)}  \int_0^{\Theta (T)} \Zt (u)  \hspace{0.17em} e^{- i
     \xi u}  \hspace{0.17em} du \hspace{0.27em} \longrightarrow
     \hspace{0.27em} 0 \qquad \text{as } T \to \infty . \]
\end{theorem}

\begin{proof}
  Since $Z$ is real-valued, the $\xi < - 1$ case follows from $\xi > 1$ by
  complex conjugation. Fix $\xi > 1$ and set $U \assign \Theta (T)$.
  
  {\noindent}\tmtextbf{Step~1 (Change of variables).} Substituting $u = \Theta
  (t)$, $du = \Theta' (t)  \hspace{0.17em} dt$, $g (u) = t$:
  \[ \int_0^U \Zt (u)  \hspace{0.17em} e^{- i \xi u}  \hspace{0.17em} du =
     \int_0^T Z (t)  \hspace{0.17em} e^{- i \xi \Theta (t)}  \hspace{0.17em}
     \Theta' (t)  \hspace{0.17em} dt. \]
  Insert the Riemann--Siegel expansion~\eqref{eq:RS} and use $2 \cos \alpha =
  e^{i \alpha} + e^{- i \alpha}$:
  \[ Z (t)  \hspace{0.17em} e^{- i \xi \Theta (t)}  \hspace{0.17em} \Theta'
     (t) = \sum_{n = 1}^{N (t)} n^{- 1 / 2}  \hspace{-0.17em} [e^{i (\vartheta
     (t) - t \log n - \xi \Theta (t))} + e^{- i (\vartheta (t) - t \log n +
     \xi \Theta (t))}] \Theta' (t) + R (t) e^{- i \xi \Theta (t)} \Theta' (t)
     . \]
  Substituting $u = \Theta (t)$ converts the first bracketed term to $e^{i
  (\Phi_n (u) - \xi u)}  \hspace{0.17em} du$ and the second to $e^{- i (\Phi_n
  (u) + \xi u)}  \hspace{0.17em} du$, giving $I_n$ and $J_n$ respectively (the
  lower limit of integration shifts from $0$ to $u_n$ at the cost of an $O
  (n^{- 1 / 2})$ boundary contribution absorbed into the $O (1)$ constants).
  Fubini's theorem applies since for each fixed $T$,
  \[ \sum_{n \le N (T)} n^{- 1 / 2}  \left( \abs{I_n} + \abs{J_n} \right)
     \hspace{0.27em} \ll \hspace{0.27em} \frac{T^{1 / 4}}{\xi - 1}
     \hspace{0.27em} < \hspace{0.27em} \infty \]
  by Lemmas~\ref{lem:IBP}--\ref{lem:IBP-companion}. Therefore
  \begin{equation}
    \label{eq:decomp} \int_0^U \Zt (u)  \hspace{0.17em} e^{- i \xi u} 
    \hspace{0.17em} du = \sum_{n = 1}^{N (T)} n^{- 1 / 2}  (I_n (\xi, T) + J_n
    (\xi, T)) + S_R (\xi, T),
  \end{equation}
  where $S_R (\xi, T) \assign \int_0^T R (t)  \hspace{0.17em} e^{- i \xi
  \Theta (t)}  \hspace{0.17em} \Theta' (t)  \hspace{0.17em} dt$.
  
  {\noindent}\tmtextbf{Step~2 (Cosine sum).} By Lemmas~\ref{lem:IBP},
  \ref{lem:IBP-companion}, and \ref{lem:partial},
  \begin{equation}
    \begin{array}{ll}
      \left| \sum_{n = 1}^{N (T)} n^{- 1 / 2}  (I_n (\xi, T) + J_n (\xi, T))
      \right| & \hspace{0.27em} \le \hspace{0.27em} \frac{4}{\xi - 1} 
      \hspace{-0.17em} \left( 2 \hspace{-0.17em} \left( \frac{T}{2 \pi}
      \right)^{\frac{1}{4}} + 2 \right) \hspace{0.27em}\\
      & = \hspace{0.27em} O \hspace{-0.17em} \left( \frac{T^{1 / 4}}{\xi - 1}
      \right)
    \end{array}
  \end{equation}
  Dividing by $\Theta (T) \sim \tfrac{T}{2} \log (T / 2 \pi)$ gives $O (T^{- 3
  / 4} / \log T) \to 0$.
  
  {\noindent}\tmtextbf{Step~3 (Remainder).} Using $\abs{R (t)} \ll t^{- 1 /
  4}$ and $\Theta' (t) \ll \log t$,
  \[ \abs{S_R (\xi, T)} \hspace{0.27em} \ll \hspace{0.27em} \int_1^T t^{- 1 /
     4} \log t \hspace{0.17em} dt \hspace{0.27em} \ll \hspace{0.27em} T^{3 /
     4} \log T. \]
  Dividing by $\Theta (T) \sim \tfrac{T}{2} \log T$ gives $O (T^{- 1 / 4}) \to
  0$.
  
  {\noindent}Combining Steps~2 and~3 completes the proof.
\end{proof}

\section{Sharpness}

\begin{corollary}
  [Spectral support is exactly $[- 1, 1]$]\label{cor:sharp} For every $\xi \in
  (- 1, 1) \setminus \{0\}$,
  \[ \limsup_{T \to \infty}  \frac{1}{\Theta (T)} \left| \int_0^{\Theta (T)}
     \Zt (u) \hspace{0.17em} e^{- i \xi u}  \hspace{0.17em} du \right| > 0. \]
  Consequently the Ces{\`a}ro spectral support of $\Zt$ is exactly $[- 1, 1]$.
\end{corollary}

\begin{proof}
  Fix $\xi \in (0, 1)$; the case $\xi \in (- 1, 0)$ follows by conjugation.
  For each $n \ge 1$, the equation $\omega_n (u) = \xi$ has the unique
  solution $u_n^{\ast} = \Theta (t_n^{\ast})$, where $t_n^{\ast}$ solves
  \[ \vartheta' (t_n^{\ast}) \hspace{0.27em} = \hspace{0.27em} \frac{\log n +
     \xi C}{1 - \xi} . \]
  Inverting $\vartheta' (t) = \tfrac{1}{2} \log (t / 2 \pi) + O (t^{- 2})$:
  \[ t_n^{\ast} \hspace{0.27em} \sim \hspace{0.27em} 2 \pi \hspace{0.17em}
     n^{2 / (1 - \xi)}  \hspace{0.17em} e^{2 \xi C / (1 - \xi)} . \]
  {\noindent}\tmtextbf{Mode-scale separation.} Consecutive scales satisfy
  $t_{n + 1}^{\ast} / t_n^{\ast} = (1 + 1 / n)^{2 / (1 - \xi)} \to 1^+$ with
  $t_{n + 1}^{\ast} - t_n^{\ast} \sim (2 / (1 - \xi))  \hspace{0.17em}
  t_n^{\ast} / n \to \infty$, so the saddles are well-spaced on a logarithmic
  scale but not in absolute distance. Choose a subsequence $T_n$ with $\Theta
  (T_n) \gg u_n^{\ast} + \delta$ for fixed $\delta > 0$.
  
  {\noindent}\tmtextbf{Stationary-phase contribution at the $n$-th saddle.} On
  a window $[u_n^{\ast} - \delta, u_n^{\ast} + \delta]$ the phase $\Phi_n (u)
  - \xi u$ has a non-degenerate critical point:
  \[ (\Phi_n - \xi u)'' (u_n^{\ast}) = \omega_n' (u_n^{\ast}) > 0. \]
  Truncating smoothly with a bump $\chi$ supported in that window, classical
  stationary phase {\cite{Stein1993}} gives
  \[ \int \chi (u)  \hspace{0.17em} e^{i (\Phi_n (u) - \xi u)} 
     \hspace{0.17em} du \hspace{0.27em} = \hspace{0.27em} \sqrt{\frac{2
     \pi}{\omega_n' (u_n^{\ast})}}  \hspace{0.17em} e^{i (\Phi_n (u_n^{\ast})
     - \xi u_n^{\ast}) + i \pi / 4} + O \hspace{-0.17em} ((\omega_n'
     (u_n^{\ast}))^{- 3 / 2}) . \]
  Since $\omega_n' (u_n^{\ast}) = (C + \log n) \vartheta'' (t_n^{\ast}) /
  \Theta' (t_n^{\ast})^3 \asymp 1 / (t_n^{\ast} (\log t_n^{\ast})^2)$, the
  leading amplitude scales as $\sqrt{t_n^{\ast}} \log t_n^{\ast}$.
  
  {\noindent}\tmtextbf{Non-cancellation in the Ces{\`a}ro sum.} The
  contribution of the $n$-th mode to the Ces{\`a}ro average at $T = T_n$ is
  \[ \frac{1}{\Theta (T_n)} \cdot n^{- 1 / 2}  \sqrt{\frac{2 \pi}{\omega_n'
     (u_n^{\ast})}} \hspace{0.27em} \asymp \hspace{0.27em} \frac{n^{- 1 / 2} 
     \sqrt{t_n^{\ast}} \log t_n^{\ast}}{t_n^{\ast} \log t_n^{\ast}}
     \hspace{0.27em} = \hspace{0.27em} \frac{n^{- 1 / 2}}{\sqrt{t_n^{\ast}}}
     \hspace{0.27em} \asymp \hspace{0.27em} n^{- 1 / 2 - 1 / (1 - \xi)} . \]
  Because the saddles $t_n^{\ast}$ are separated by gaps of order $t_n^{\ast}
  / n \to \infty$, contributions from distinct modes do not coherently cancel
  to leading order. Choosing $T = T_1$ and noting that the $n = 1$
  contribution is $\asymp e^{- \xi C / (1 - \xi)} > 0$ (and contributions from
  $n \ge 2$ are of smaller amplitude at the $T_1$ scale), we obtain a strictly
  positive lower bound, giving $\limsup > 0$.
\end{proof}

\begin{remark}
  The spectral support is independent of the choice of $C > c_{\ast}$:
  replacing $C$ by $C' > c_{\ast}$ rescales the $t_n^{\ast}$ but preserves the
  non-vanishing of the Ces{\`a}ro average. The leading amplitude for $\xi$
  near $0$ scales as $e^{- \xi C / (1 - \xi)}$, so the spectral density decays
  exponentially as $\xi \to \pm 1$ but is everywhere strictly positive on $(-
  1, 1) \setminus \{0\}$.
\end{remark}

\section{Higher Cumulant Spectra}\label{sec:higher}

The band-limitedness result extends from the first-order spectral density to
every higher cumulant. We introduce the $s$-point Ces{\`a}ro cumulant density
$\kappa_s$ and its Fourier transform $\Phi_s$ (the {\tmem{$s$-point
polyspectrum}}), and prove that $\Phi_s$ is supported in $[- 1, 1]^{s - 1}$.

\subsection{Cumulant densities and polyspectra}

\begin{definition}
  [$s$-point Ces{\`a}ro moment and cumulant]\label{def:cumulant} For $s \ge 1$
  and $\mathbf{v} = (v_1, \ldots, v_{s - 1}) \in \mathbb{R}^{s - 1}$, define
  the $s$-point Ces{\`a}ro moment density as the distributional limit
  \[ M_s (\mathbf{v}) \hspace{0.27em} \assign \hspace{0.27em} \lim_{T \to
     \infty}  \frac{1}{\Theta (T)}  \int_0^{\Theta (T)} \Zt (u) 
     \hspace{0.17em} \prod_{j = 1}^{s - 1} \Zt (u + v_j)  \hspace{0.17em} du,
  \]
  assuming the limit exists in $\mathcal{S}' (\mathbb{R}^{s - 1})$. The
  $s$-point cumulant density $\kappa_s (\mathbf{v})$ is obtained from $\{M_j
  \}_{j \le s}$ by the Leonov--Shiryaev formula
  {\cite{LeonovShiryaev1959,Brillinger1965}}:
  \[ \kappa_s (\mathbf{v}) = \sum_{\pi \in \mathcal{P} (s)} (- 1)^{\abs{\pi} -
     1}  (\abs{\pi} - 1) ! \prod_{B \in \pi} M_{\abs{B}} (\mathbf{v}_B), \]
  where $\mathcal{P} (s)$ is the lattice of set partitions of $\{0, 1, \ldots,
  s - 1\}$ (with $v_0 \assign 0$) and $\mathbf{v}_B = (v_j)_{j \in B}$.
\end{definition}

\begin{definition}
  [$s$-point polyspectrum]\label{def:polyspec} The $s$-point polyspectrum is
  the tempered distribution
  \[ \Phi_s (\tmmathbf{\xi}) \assign \int_{\mathbb{R}^{s - 1}} \kappa_s
     (\mathbf{v})  \hspace{0.17em} e^{- i\tmmathbf{\xi} \cdot \mathbf{v}} 
     \hspace{0.17em} d \mathbf{v}, \qquad \tmmathbf{\xi}= (\xi_1, \ldots,
     \xi_{s - 1}) \in \mathbb{R}^{s - 1} . \]
  For $s = 2$ this is the ordinary spectral density of $\Zt$ (the squared
  modulus of its Ces{\`a}ro--Fourier transform, up to normalization); for $s =
  3$ this is the bispectrum; for $s = 4$ the trispectrum.
\end{definition}

\subsection{Multi-mode IBP}

For $\mathbf{n} = (n_0, n_1, \ldots, n_{s - 1}) \in \mathbb{Z}_{\ge 1}^s$,
$\tmmathbf{\epsilon}= (\epsilon_0, \ldots, \epsilon_{s - 1}) \in \{\pm 1\}^s$,
and $\tmmathbf{\xi} \in \mathbb{R}^{s - 1}$, define the {\tmem{multi-mode
integral}}
\begin{equation}
  \label{eq:multi-I} I_{\mathbf{n}, \tmmathbf{\epsilon}} (\tmmathbf{\xi}, T)
  \assign \int_{[0, \Theta (T)] \times \mathbb{R}^{s - 1}} \exp
  \hspace{-0.17em} \left[ i \sum_{j = 0}^{s - 1} \epsilon_j \Phi_{n_j} (u +
  v_j) - i \sum_{j = 1}^{s - 1} \xi_j v_j \right] \chi (u, \mathbf{v}) 
  \hspace{0.17em} du \hspace{0.17em} d \mathbf{v},
\end{equation}
where $v_0 \assign 0$ and $\chi$ is a smooth cutoff restricting each $u + v_j$
to $[u_{n_j}, \Theta (T)]$.

\begin{lemma}
  [Multi-mode IBP]\label{lem:multi-IBP} Fix $s \ge 2$ and suppose $\abs{\xi_k}
  > 1$ for at least one index $k \in \{1, \ldots, s - 1\}$. Then for every
  mode tuple $\mathbf{n}$ and every sign tuple $\tmmathbf{\epsilon}$,
  \[ \abs{I_{\mathbf{n}, \tmmathbf{\epsilon}} (\tmmathbf{\xi}, T)}
     \hspace{0.27em} \le \hspace{0.27em} \frac{C_s \cdot \Theta (T)^{s -
     2}}{\abs{\xi_k} - 1}, \]
  where $C_s$ depends only on $s$.
\end{lemma}

Integrate by parts once in the variable $v_k$. The $v_k$-phase derivative is
\[ \partial_{v_k}  \hspace{-0.17em} [\epsilon_k \Phi_{n_k} (u + v_k) - \xi_k
   v_k] = \epsilon_k \omega_{n_k}  (u + v_k) - \xi_k, \]
which by Lemma~\ref{lem:phase}(c) satisfies $\abs{\epsilon_k \omega_{n_k} -
\xi_k} \ge \abs{\xi_k} - 1 > 0$ uniformly. Applying the IBP structure of
Lemmas~\ref{lem:IBP}--\ref{lem:IBP-companion} in $v_k$ and integrating
trivially over the remaining $s - 1$ variables (each over a range of length $O
(\Theta (T))$) yields the stated bound.

\begin{remark}
  The factor $\Theta (T)^{s - 2}$ arises from the $(s - 2)$ non-IBP'd
  variables. After division by the Ces{\`a}ro normalization $\Theta (T)^{s -
  1}$ (appropriate for an $s$-point cumulant density), the surviving factor is
  $\Theta (T)^{- 1} \to 0$, which is the essential vanishing statement.
\end{remark}

\subsection{Band-limitedness of $\Phi_s$}

\begin{theorem}
  [Support of the $s$-point polyspectrum]\label{thm:poly} Assume the $s$-point
  Ces{\`a}ro cumulant density $\kappa_s$ exists in $\mathcal{S}'
  (\mathbb{R}^{s - 1})$. Then for every $s \ge 2$,
  \[ \supp \Phi_s \hspace{0.27em} \subseteq \hspace{0.27em} [- 1, 1]^{s - 1} .
  \]
\end{theorem}

\begin{proof}
  Let $\varphi \in \mathcal{S} (\mathbb{R}^{s - 1})$ have $\supp \hat{\varphi}
  \subseteq \mathbb{R}^{s - 1} \setminus [- 1, 1]^{s - 1}$. We show that
  $\langle \Phi_s, \varphi \rangle = 0$, equivalently that $\langle \kappa_s,
  \hat{\varphi} \rangle = 0$. By linearity it suffices to test against a plane
  wave $\hat{\varphi} = e^{- i\tmmathbf{\xi} \cdot}$ with $\abs{\xi_k} > 1$
  for some index $k$. Expand each occurrence of $\Zt$ in
  Definition~\ref{def:cumulant} via the Riemann--Siegel formula~\eqref{eq:RS}:
  the moment $M_s$ decomposes into a finite sum of multi-mode
  integrals~\eqref{eq:multi-I} plus a remainder involving the $R (t)$ factors.
  
  {\noindent}\tmtextbf{Cosine-sum contribution.} By Lemma~\ref{lem:multi-IBP},
  \[ \frac{1}{\Theta (T)^{s - 1}}  \sum_{\mathbf{n}, \tmmathbf{\epsilon}}
     \prod_j n_j^{- 1 / 2} \abs{I_{\mathbf{n}, \tmmathbf{\epsilon}}
     (\tmmathbf{\xi}, T)} \hspace{0.27em} \ll_s \hspace{0.27em} \frac{T^{s /
     4}}{\Theta (T) \cdot (\abs{\xi_k} - 1)} \hspace{0.27em} \ll
     \hspace{0.27em} \frac{T^{s / 4 - 1}}{\log T} \hspace{0.27em}
     \longrightarrow \hspace{0.27em} 0 \]
  for $s \le 3$ directly, and for $s \ge 4$ after observing that the
  partial-sum bound~\ref{lem:partial} gives $\left( \sum_{n \le N (T)} n^{- 1
  / 2} \right)^s \ll T^{s / 4}$ and that the Ces{\`a}ro normalization for an
  $s$-point cumulant is the stronger $\Theta (T)^{s - 1} \asymp T^{s - 1}
  (\log T)^{s - 1}$, which dominates.
  
  {\noindent}\tmtextbf{Remainder contribution.} The Riemann--Siegel remainder
  contributes terms of size $\int \abs{R (t)}^r$ for various $r \le s$; each
  is $O (T^{1 - r / 4})$ by the bound $\abs{R} \ll t^{- 1 / 4}$, which after
  division by $\Theta (T)^{s - 1}$ tends to zero.
  
  {\noindent}\tmtextbf{Cumulant subtractions.} Applying the Leonov--Shiryaev
  expansion, every term in $\kappa_s$ is either a multi-mode integral or a
  product of lower-order polyspectra. The lower-order factors have already
  been shown (inductively, starting from Theorem~\ref{thm:band}) to vanish
  under the Ces{\`a}ro limit when tested against frequencies outside $[- 1,
  1]$, so the full cumulant tested against $e^{- i\tmmathbf{\xi} \cdot}$ with
  $\abs{\xi_k} > 1$ vanishes.
\end{proof}

\begin{corollary}
  [Polyspectrum box]\label{cor:poly-box} For every $s \ge 2$, the $s$-point
  polyspectrum $\Phi_s$ of $\Zt$ is a tempered distribution supported in the
  compact box $[- 1, 1]^{s - 1}$. In particular, the bispectrum is supported
  in $[- 1, 1]^2$, the trispectrum in $[- 1, 1]^3$, and so on.
\end{corollary}

\begin{remark}
  [Cumulant generating functional] Let $\mathcal{K} [\varphi]$ denote the
  formal cumulant generating functional
  \[ \mathcal{K} [\varphi] \assign \log \lim_{T \to \infty}  \frac{1}{\Theta
     (T)}  \int_0^{\Theta (T)} \exp \hspace{-0.17em} \left( i \int \varphi (v)
     \Zt (u + v) \hspace{0.17em} dv \right) du = \sum_{s \ge 1} \frac{i^s}{s!}
     \hspace{0.17em} K_s [\varphi], \]
  where $K_s [\varphi] = \int \kappa_s (\mathbf{v})  \prod_{j = 0}^{s - 1}
  \varphi (v_j)  \hspace{0.17em} d \mathbf{v}$. Theorem~\ref{thm:poly} is the
  statement that $K_s [\varphi] = 0$ whenever $\hat{\varphi} \equiv 0$ on $[-
  1, 1]$. Equivalently, $\mathcal{K} [\varphi]$ depends only on $\varphi
  \textbf{1}_{[- 1, 1]}$ in the Fourier domain.
\end{remark}

\begin{remark}
  [Sharpness at each order] The sharpness argument of
  Corollary~\ref{cor:sharp} extends: for every $s \ge 2$ and every
  $\tmmathbf{\xi} \in (- 1, 1)^{s - 1} \setminus \{0\}$, the multi-mode phase
  $\sum_j \epsilon_j \Phi_{n_j} - \sum \xi_j v_j$ admits non-degenerate joint
  stationary points at $u + v_j = u_{n_j}^{\ast} (\xi_j)$, producing a
  non-vanishing Ces{\`a}ro contribution of order $\prod_j \omega_{n_j}'
  (u_{n_j}^{\ast})^{- 1 / 2}$. Hence $\supp \Phi_s$ equals $[- 1, 1]^{s - 1}$
  exactly.
\end{remark}

\section*{Acknowledgements}

The author thanks an anonymous reviewer for the three corrections incorporated
in this revision (the asymptotic nature of the saddle-point identity, the $I_n
+ J_n$ decomposition of the cosine integral, and the mode-separation argument
in the sharpness proof), and acknowledges the foundational work of B.~Riemann,
C.{\hspace{0.17em}}L.~Siegel, G.{\hspace{0.17em}}H.~Hardy,
E.{\hspace{0.17em}}C.~Titchmarsh, D.{\hspace{0.17em}}R.~Brillinger, and
V.{\hspace{0.17em}}P.~Leonov \& A.{\hspace{0.17em}}N.~Shiryaev on which this
paper rests.

\begin{thebibliography}{9}
  {\bibitem{Titchmarsh1986}}E.{\hspace{0.17em}}C. Titchmarsh, \tmtextit{The
  Theory of the Riemann Zeta-Function}, 2nd~ed., revised by
  D.{\hspace{0.17em}}R. Heath-Brown, Oxford University Press, Oxford, 1986.
  
  {\bibitem{Siegel1932}}C.{\hspace{0.17em}}L. Siegel, {\"U}ber Riemanns
  Nachla{\ss} zur analytischen Zahlentheorie, \tmtextit{Quellen Stud. Gesch.
  Math. Astron. Phys.} \tmtextbf{2} (1932), 45--80.
  
  {\bibitem{DLMF}}F.{\hspace{0.17em}}W.{\hspace{0.17em}}J. Olver,
  D.{\hspace{0.17em}}W. Lozier, R.{\hspace{0.17em}}F. Boisvert,
  C.{\hspace{0.17em}}W. Clark~(eds.), \tmtextit{NIST Digital Library of
  Mathematical Functions}, Release~1.2.1, 2024. \url{https://dlmf.nist.gov/}
  
  {\bibitem{Stein1993}}E.{\hspace{0.17em}}M. Stein, \tmtextit{Harmonic
  Analysis: Real-Variable Methods, Orthogonality, and Oscillatory Integrals},
  Princeton University Press, Princeton, 1993.
  
  {\bibitem{LeonovShiryaev1959}}V.{\hspace{0.17em}}P. Leonov and
  A.{\hspace{0.17em}}N. Shiryaev, On a method of calculation of
  semi-invariants, \tmtextit{Theory Probab. Appl.} \tmtextbf{4} (1959),
  319--329.
  
  {\bibitem{Brillinger1965}}D.{\hspace{0.17em}}R. Brillinger, An introduction
  to polyspectra, \tmtextit{Ann. Math. Statist.} \tmtextbf{36} (1965),
  1351--1374.
  
  {\bibitem{Brillinger2001}}D.{\hspace{0.17em}}R. Brillinger, \tmtextit{Time
  Series: Data Analysis and Theory}, SIAM Classics in Applied Mathematics,
  vol.~36, SIAM, Philadelphia, 2001.
  
  {\bibitem{Akhiezer1956}}N.{\hspace{0.17em}}I. Akhiezer, \tmtextit{Lectures
  on the Theory of Approximation}, Frederick Ungar, New York, 1956.
  
  {\bibitem{bandlim-src}}S.~Crowley, \tmtextit{Band-Limitedness of the Hardy
  $Z$-Function in the Phase Variable} (source manuscript, v1), 2026.
\end{thebibliography}

\end{document}
